\chapter*{Ringraziamenti}
\thispagestyle{empty}

Vorrei innanzitutto ringraziare il mio relatore, il Prof. Lorenzo Bruzzone, e il mio correlatore,
il Dott. Gianmarco Perantoni, per avermi dato l'opportunità di svolgere questo lavoro di tesi presso
il loro laboratorio e per avermi seguito e guidato con competenza durante l'intero percorso di ricerca.\\

Un ringraziamento speciale va alla mia famiglia e a tutti i miei parenti, per avermi supportato e
incoraggiato in questi anni di studio, permettendomi di raggiungere questo importante traguardo.\\

Ai miei compagni di corso, che mi hanno accompagnato in questi tre anni intensi, condividendo gioie e difficoltà.
Con loro è nato un saldo rapporto di amicizia e di aiuto reciproco che spero possa continuare ben oltre la fine di questo percorso.\\

Grazie a tutti i miei amici, vicini e lontani, storici e recenti: siete sempre stati un punto di
riferimento fondamentale e sono ogni giorno più fiero di essere vostro amico.\\

Un pensiero va anche agli amici dello studentato San Bartolomeo, presenti o passati, e a tutte le
persone che hanno contribuito a rendere unica e formativa la mia esperienza universitaria.\\

Infine, voglio dedicare un piccolo ringraziamento e una citazione a Caparezza,
che con la sua canzone \textit{Curiosity} ha acceso in me l'interesse per l'osservazione della Terra e per l'argomento di questo elaborato:

\begin{quote}
  \textit{
    ''Con lo shuttle o con la stazione spaziale si va alle stesse quote, praticamente,
    tra i trecentottanta e i quattrocento chilometri di quota.
    Si va a circa tre ore di macchina se uno potesse guidare verso l'alto e quello che
    colpisce è che si vedono le città e come cambia la morfologia del terreno a causa dell'uomo, ma non si vede l'uomo.''
  }
\end{quote}

E per concludere, prendo in prestito ancora una volta le sue parole, che riflettono appieno il mio profondo rispetto per l'istruzione e la scoperta:

\begin{quote}
  \textit{
    ''Siamo tutti qui per imparare, non per fare i conti in tasca agli astrofisici mentre gettiamo spicci nelle fontane.
    La curiosità non è un istinto, è un dovere morale, è un dovere morale.''
  }
\end{quote}

In un mondo in cui la curiosità e l'istruzione vengono sempre più spesso messe da parte,
è fondamentale ricordare che la conoscenza è un diritto e un dovere di tutti,
e che la ricerca scientifica è il motore del progresso umano.